% theory/01-memoria.tex — T06: Part 1, the price of memory
% Project SemComV1 · feature 001-codigo-trayectoria · 2026-10-13
% Notation and definitions: theory/00-modelo.tex. Numbers: ejemplo_semana.py, contraejemplo_memoria.py.
% Numerical checks of every statement: check_t06.py, check_t06b.py (constitution P4.4).
\documentclass[11pt]{article}
\usepackage[margin=1in]{geometry}
\usepackage{amsmath,amssymb,amsthm,mathtools}
\usepackage{booktabs}
\usepackage{tikz}
\usetikzlibrary{arrows.meta,positioning}
\usepackage[hidelinks]{hyperref}
\newcommand{\doi}[1]{doi:\,\url{#1}}

\newtheorem{theorem}{Theorem}
\newtheorem{proposition}{Proposition}
\newtheorem{corollary}{Corollary}
\newtheorem{lemma}{Lemma}
\theoremstyle{definition}
\newtheorem{definition}{Definition}
\theoremstyle{remark}
\newtheorem{remark}{Remark}
\newtheorem{example}{Example}

\newcommand{\M}{\mathcal{M}}
\newcommand{\Pset}{\Pi}
\newcommand{\E}{\mathbb{E}}
\newcommand{\D}{D}
\DeclareMathOperator*{\argmax}{arg\,max}

\title{Path Codes on Shared Knowledge Graphs\\[2pt]\large Part 1: The Price of Memory}
\author{R.~Sandoval Ar\'echiga (working notes)}
\date{October 13, 2026 --- v0.1}

\begin{document}
\maketitle

\begin{abstract}
A path code conveys a meaning by describing a walk on a shared graph, decision by decision. We compare three ways of doing so. \emph{Native} codes use only the shared routing statistics and pay, above the semantic entropy $H(M)$, the average conditional min-entropy of the walk given the meaning. \emph{Adapted} codes, which also share the selection of one walk per meaning, attain $H(M)$. They need no decoder memory when the meaning is determined by the endpoint of the walk, but may need it when the meaning depends on the route (context-dependent meaning, or polysemy). \emph{Bits-back} codes attain $H(M)$ with native statistics, at the price of coupling consecutive messages. A corollary quantifies when the semantic Markov sources of Xu et al.~\cite{xu2025} have the entropy their formula assigns to them.
\end{abstract}

\section{Setting}

We use the model of Part~0: a trajectory network $G$ satisfying Assumptions~1--3, path measure $P$, meaning map $\mu:\Pset_G\to\M$, preimages $\Pset(m)$, and meaning law $P_M$. Code lengths are \emph{ideal}: a sequential code that uses distribution $Q$ assigns $-\log_2 Q(\pi)$ bits to walk $\pi$. By Lemma~1 of Part~0, arithmetic coding realises any such length within $2$ bits~\cite{rissanen1979,witten1987}, so every statement below holds for real codes up to an additive $2$.

\begin{definition}[Selection and induced law]
A \emph{selection} is a map $f:\M\to\Pset_G$ with $f(m)\in\Pset(m)$. It induces the law $Q_f$ on walks given by $Q_f(f(m))=P_M(m)$ and $Q_f(\pi)=0$ for walks outside the image of $f$.
\end{definition}

\begin{definition}[Native, adapted and bits-back path codes]\label{def:codes}
\begin{enumerate}
\item A \emph{native} code with selection $f$ describes $f(m)$ using the shared routing kernel, at cost $-\log_2 P(f(m))$. Its rate is $L_{\mathrm{nat}}(f)=\E[-\log_2 P(f(M))]$. Transmitter and receiver share only $G$ and $\mu$.
\item An \emph{adapted} code with selection $f$ and state space $S$ is a finite-state decoder. When the walk so far is $\rho$ and ends at node $v$, the decoder is in state $s(\rho)\in S$. It decodes the next decision with a kernel $q(\cdot\mid v,s)$, and moves to $s(\rho e)=\delta(s(\rho),e)$. Its rate is $L_{\mathrm{ad}}=\E_{Q_f}[-\log_2 q(\Pi)]$, where $q(\pi)$ is the product of the kernels along $\pi$. It is \emph{memoryless} if $|S|=1$. In addition to $G$ and $\mu$, both ends share $f$, $q$ and $\delta$.
\item A \emph{bits-back} code uses the native kernel and a randomised selection $\pi\sim P(\cdot\mid m)$, and recovers the selection bits at the receiver~\cite{hinton1993,frey1997}.
\end{enumerate}
\end{definition}

Every path code lets the receiver recover the meaning, so its rate is at least $H(M)$, the semantic entropy $H_s$ of~\cite{niu2025}.

\section{Native codes}

For $m\in\M$ let $\pi^\star(m)\in\argmax_{\pi\in\Pset(m)}P(\pi)$, with ties broken by a fixed order, and define
\begin{equation}\label{eq:g}
g(m)=\log_2\frac{P_M(m)}{P(\pi^\star(m))}=-\log_2\max_{\pi\in\Pset(m)}P(\pi\mid M=m).
\end{equation}
Thus $g(m)$ is the min-entropy of the walk given the meaning, $H_\infty(\Pi\mid M=m)$.

\begin{theorem}[Price of native routing]\label{thm:native}
\begin{enumerate}
\item[(a)] $\min_f L_{\mathrm{nat}}(f)=L_{\mathrm{nat}}(\pi^\star)=H(M)+\E[g(M)]$.
\item[(b)] For every $m$, $0\le g(m)\le H(\Pi\mid M=m)\le\log_2|\Pset(m)|$. Hence $0\le\E[g(M)]\le H(\Pi\mid M)$.
\item[(c)] $\E[g(M)]=0$ if and only if every meaning has exactly one walk of positive probability. $\E[g(M)]=H(\Pi\mid M)$ if and only if, for every $m$ with $P_M(m)>0$, the conditional law $P(\cdot\mid M=m)$ is uniform on its support.
\end{enumerate}
\end{theorem}

\begin{proof}
(a) The cost of a native code is $\sum_m P_M(m)\,(-\log_2 P(f(m)))$, which is minimised term by term by $f(m)=\pi^\star(m)$. Writing $-\log_2 P(\pi^\star(m))=-\log_2 P_M(m)+g(m)$ and averaging gives the claim.

(b) $g(m)\ge0$ because $P(\pi^\star(m))\le P_M(m)$. Fix $m$ and write $c(\pi)=P(\pi\mid M=m)$. Then $g(m)=-\log_2\max_\pi c(\pi)\le\sum_\pi c(\pi)(-\log_2 c(\pi))=H(\Pi\mid M=m)$, because $-\log_2\max_\pi c(\pi)\le-\log_2 c(\pi)$ for every $\pi$. The last inequality is the usual bound of entropy by the logarithm of the support size~\cite{cover2006}. Averaging over $M$ gives the second statement.

(c) $\E[g(M)]=0$ if and only if $g(m)=0$ for every $m$ with $P_M(m)>0$, that is, $\max_\pi c(\pi)=1$. For the upper case, $\E[g(M)]=H(\Pi\mid M)$ requires equality in (b) for every such $m$. The inequality $-\log_2\max c\le-\log_2 c(\pi)$ is tight in average if and only if $c(\pi)=\max c$ on the support, that is, $c$ is uniform there.
\end{proof}

\begin{remark}[Kraft slack]
Native lengths waste coding space:
\[
\sum_m 2^{\log_2P(\pi^\star(m))}=\sum_mP(\pi^\star(m))=1-\sum_m\sum_{\pi\in\Pset(m)\setminus\{\pi^\star(m)\}}P(\pi).
\] The mass of the unused synonymous walks is coding space that no message ever uses.
\end{remark}

\section{Adapted codes and decoder memory}

For a selection $f$ and a prefix $\rho$ of some walk in the image of $f$, let $Q_f(\cdot\mid\rho)$ be the law of the remaining decisions, given that the selected walk starts with $\rho$. For a node $v$, let
\[
N_v(f)=\bigl|\{Q_f(\cdot\mid\rho):\rho\text{ ends at }v,\ Q_f(\rho)>0\}\bigr|
\]
be the number of distinct future laws at $v$, and $N(f)=\max_vN_v(f)$.

\begin{theorem}[Adapted codes]\label{thm:adapted}
\begin{enumerate}
\item[(a)] For every selection $f$ there is an adapted code with $L_{\mathrm{ad}}=H(M)$.
\item[(b)] An adapted code with selection $f$ attains $L_{\mathrm{ad}}=H(M)$ if and only if its state assignment separates prefixes that end at the same node and have different future laws. The least number of states that suffices is $N(f)$. Any code that merges two such prefixes has $L_{\mathrm{ad}}>H(M)$.
\item[(c)] If the meaning is \emph{endpoint-determined} ($\mu(\pi)$ depends only on the last node of $\pi$), then the Viterbi selection $\pi^\star$ satisfies $N(\pi^\star)=1$, so a memoryless adapted code attains $H(M)$.
\item[(d)] There exist networks with route-dependent meaning for which $N(f)\ge2$ for every selection $f$. For them, every memoryless adapted code has $L_{\mathrm{ad}}>H(M)$.
\end{enumerate}
\end{theorem}

\begin{proof}
(a) Since $f$ is injective, $Q_f$ is $P_M$ transported to the image of $f$, so $H(Q_f)=H(M)$. Take $S$ to be the set of prefixes and $q(\cdot\mid v,\rho)=Q_f(\cdot\mid\rho)$. The kernels then multiply to $Q_f(\pi)$ along each selected walk, and the rate is $\E_{Q_f}[-\log_2 Q_f(\Pi)]=H(Q_f)=H(M)$.

(b) For any adapted code, $L_{\mathrm{ad}}=\E_{Q_f}[-\log_2 q(\Pi)]=H(Q_f)+\D(Q_f\|q)$, where $q$ denotes the walk law the code induces. By the chain rule for relative entropy,
\[
\D(Q_f\|q)=\sum_{\rho}Q_f(\rho)\,\D\bigl(Q_f(\cdot\mid\rho)_1\,\big\|\,q(\cdot\mid v(\rho),s(\rho))\bigr),
\]
where the sum runs over prefixes, $Q_f(\cdot\mid\rho)_1$ is the law of the next decision, and $v(\rho)$ is the last node of $\rho$. The rate equals $H(M)$ if and only if every term vanishes, i.e.\ the kernel at $(v(\rho),s(\rho))$ equals the next-decision law of every prefix assigned to that state. Two prefixes ending at $v$ with different future laws differ in the next-decision law at some later prefix along a common continuation, so a state map that merges them makes some term positive. Conversely, the map ``prefix $\mapsto$ its future law'' is a valid state assignment: the future law of $\rho e$ is the future law of $\rho$ conditioned on the first decision being $e$, hence a function of the future law of $\rho$ and of $e$. This defines $\delta$, and uses exactly $N_v(f)$ states at $v$.

(c) Define $\rho^\star(r)=(r)$ and, in a topological order of the DAG, $\rho^\star(v)=\rho^\star(u)\,v$ with $u$ a predecessor maximising $P(\rho^\star(u))p(v\mid u)$, ties broken by a fixed order. By induction, $\rho^\star(v)$ is a most probable walk from $r$ to $v$, and the set $\{\rho^\star(v)\}_v$ is prefix-closed, i.e.\ it forms a tree rooted at $r$. Since the meaning is endpoint-determined, $\max_{\pi\in\Pset(m)}P(\pi)=\max_{t:\,\mu(t)=m}P(\rho^\star(t))$. Choosing for each $m$ a maximising terminal $t_m$ and setting $\pi^\star(m)=\rho^\star(t_m)$ gives a Viterbi selection whose walks are branches of that tree. Every node is then reached by at most one selected prefix, so $N_v(\pi^\star)\le1$ for every $v$.

(d) Example~\ref{ex:poly} below has $N(f)=2$ for each of its selections. The strict inequality then follows from (b).
\end{proof}

\begin{remark}[What memory buys]
Theorems~\ref{thm:native} and~\ref{thm:adapted} separate three resources. The native code needs only the shared graph and pays $\E[g(M)]$. The adapted code removes that payment but requires sharing a selection; it never requires decoder memory when the meaning is endpoint-determined, and it can require it when the meaning of a decision depends on how the walk arrived there. In the language of the original note, \emph{polysemy forces the walker to carry state}: Theorem~\ref{thm:adapted}(c)--(d) makes this precise in one direction and shows by example that the other can fail. The least memory $\min_fN(f)$ is a combinatorial quantity whose computational complexity is left open.
\end{remark}

\section{Bits-back codes}

\begin{proposition}[Native statistics, amortised optimality]\label{prop:bb}
Under the standard idealised model of bits-back coding~\cite{hinton1993,frey1997}, the expected total length to send $n$ i.i.d.\ meanings is at most $nH(M)+H(\Pi\mid M)+2$. The rate therefore tends to $H(M)$ using only $G$ and $\mu$.
\end{proposition}

\begin{proof}[Proof sketch]
For each message the encoder draws $\pi\sim P(\cdot\mid m)$ by \emph{decoding} from a stack of previously produced bits, which consumes $-\log_2P(\pi\mid m)$ bits on average. It then encodes $\pi$ with the native kernel, which adds $-\log_2P(\pi)$ bits. The receiver decodes $\pi$, computes $m=\mu(\pi)$, and re-encodes $\pi$ under $P(\cdot\mid m)$, which restores the consumed bits. The net cost per message is $-\log_2P(\pi)+\log_2P(\pi\mid m)=-\log_2P_M(m)$, whose mean is $H(M)$. Only the first message needs fresh bits, $H(\Pi\mid M)$ on average. The final arithmetic-coding flush adds at most $2$ bits.
\end{proof}

\section{Stationary sources: when is a semantic Markov chain Markov?}

Let $(X_k)$ be a stationary Markov chain on a finite set, $\varphi$ a synonymous mapping of its states, and $Y_k=\varphi(X_k)$. Xu et al.~\cite[Def.~1]{xu2025} define the semantic transition matrix by stationary-weighted aggregation. Under stationarity this is exactly the law of $Y_2$ given $Y_1$. Their block entropy~\cite[Thm.~3]{xu2025}, for $m=1$, is
\[
\widehat H_N=H(Y_1)+(N-1)\,H(Y_2\mid Y_1).
\]

\begin{corollary}[First-order semantic models]\label{cor:xu}
\begin{enumerate}
\item[(a)] The excess of the first-order formula is
\[
\widehat H_N-H(Y_1,\dots,Y_N)=\sum_{k=3}^{N}I\bigl(Y_k;Y_1^{k-2}\mid Y_{k-1}\bigr)\ge0.
\]
\item[(b)] $\widehat H_N=H(Y_1,\dots,Y_N)$ for all $N$ if and only if $(Y_k)$ is a Markov chain. This holds, for instance, when $X$ is lumpable with respect to $\varphi$~\cite{geiger2014}.
\item[(c)] The per-symbol excess $\lim_N(\widehat H_N-H(Y_1^N))/N=H(Y_2\mid Y_1)-\bar H(Y)$ is the price of a memoryless semantic model, where $\bar H(Y)$ is the entropy rate.
\end{enumerate}
\end{corollary}

\begin{proof}
By the chain rule, $H(Y_1^N)=H(Y_1)+\sum_{k=2}^NH(Y_k\mid Y_1^{k-1})$. By stationarity, $H(Y_k\mid Y_{k-1})=H(Y_2\mid Y_1)$. Hence
\[
\widehat H_N-H(Y_1^N)=\sum_{k=2}^N\bigl[H(Y_k\mid Y_{k-1})-H(Y_k\mid Y_1^{k-1})\bigr]=\sum_{k=3}^NI(Y_k;Y_1^{k-2}\mid Y_{k-1}),
\]
where the $k=2$ term vanishes. This proves (a). Each term is zero if and only if $Y_k$ is conditionally independent of $Y_1^{k-2}$ given $Y_{k-1}$; all of them vanish if and only if $(Y_k)$ is Markov, which proves (b). For (c), divide by $N$ and use $H(Y_k\mid Y_1^{k-1})\to\bar H(Y)$.
\end{proof}

\begin{remark}
In words: the block entropy of~\cite[Thm.~3]{xu2025} is exact when the aggregated process is Markov, and in general it is the code length of a first-order model of the semantic source. In 200 random four-state chains with the partition $\{1\},\{2\},\{3,4\}$ and $N=6$, the excess reached $1.11$ bits (script \texttt{check\_xu.py}). The excess plays here the role that $\E[g(M)]$ plays for walks: it is what a decoder without memory pays.
\end{remark}

\section{Examples}

\begin{example}[Endpoint-determined meaning: Part~0 revisited]
In the ``what day is today'' network the meaning is endpoint-determined (Tuesday and Tuesday$'$ are distinct terminals). Theorem~\ref{thm:native} gives the native rate $H(M)+\E[g(M)]=1.485+0.737=2.222$ bits, with $\E[g(M)]=0.737<H(\Pi\mid M)=0.971$. The bound is not tight because $P(\cdot\mid M=m)=(0.6,0.4)$ is not uniform. By Theorem~\ref{thm:adapted}(c), the Viterbi selection routes everything through ``today is'', and a memoryless adapted kernel attains $H(M)=1.485$ bits.
\end{example}

\begin{example}[Context-dependent meaning forces memory]\label{ex:poly}
Let $r\to a$ with probability $0.6$ and $r\to b$ with probability $0.4$. Both $a$ and $b$ lead to a shared node $x$, a ``word''. From $x$ the walk ends at $t_1,t_2,t_3$ with probabilities $(0.5,0.3,0.2)$. The meaning depends on the context branch:
\begin{center}
\begin{tabular}{lcccccc}
\toprule
walk & $a\,t_1$ & $a\,t_2$ & $a\,t_3$ & $b\,t_1$ & $b\,t_2$ & $b\,t_3$\\
$P$ & $0.30$ & $0.18$ & $0.12$ & $0.20$ & $0.12$ & $0.08$\\
meaning & A & A & B & A & C & C\\
\bottomrule
\end{tabular}
\end{center}
The word $x$ is polysemous: $t_2$ means A after $a$ and C after $b$. Here $P_M=(0.68,0.12,0.20)$, $H(M)=1.210$ and $H(\Pi\mid M)=1.247$ bits.
\begin{itemize}
\item \emph{Native code:} $2.160$ bits, an excess of $0.950$ bits.
\item \emph{Any selection needs two states at $x$.} The future law at $x$ after $a$ (resp.\ $b$) is the law of the endpoints of the selected walks through $a$ (resp.\ $b$). B has the single walk $a\,t_3$, and C only walks through $b$. If A is selected through $a$, the support after $a$ has two endpoints and the support after $b$ has one. If A is selected through $b$, the support after $a$ is $\{t_3\}$ and after $b$ has two endpoints. In both cases the two laws differ, so $N(f)=2$ for all six selections ($3\times1\times2$), as an exhaustive check confirms.
\item \emph{Memoryless adapted codes:} the best one over all selections and node kernels has rate $1.251$ bits, an excess of $0.041$ bits. With the Viterbi selection the excess is $0.722$ bits.
\item \emph{Adapted code with two states at $x$:} rate $H(M)=1.210$ bits.
\end{itemize}
Script: \texttt{contraejemplo\_memoria.py}.
\end{example}

\section{What a reviewer will ask (and the answers)}

\begin{enumerate}
\item \emph{``The min-entropy bound is elementary.''} The bound is; the contribution is the separation between native, adapted and bits-back codes. That separation is not expressible in unstructured synonymous coding~\cite{niu2025,liang2025}, where a codeword names the set and no walk is traversed. It is also not found in the bits-back literature, where the latent variable is not a route on a shared graph.
\item \emph{``Adapted codes simply code the meaning.''} Yes, at rate $H(M)$. But they must be decodable decision by decision on the shared graph, and Theorem~\ref{thm:adapted}(b)--(d) shows that this constraint costs decoder memory when the meaning is context-dependent.
\item \emph{``Why not always use bits-back?''} It couples messages, needs an initial overhead of $H(\Pi\mid M)$ bits, and loses the per-message decodability that a walk provides. The three schemes trade shared knowledge, memory and latency.
\end{enumerate}

\begin{thebibliography}{9}
\bibitem{niu2025} K.~Niu and P.~Zhang, \emph{The Mathematical Theory of Semantic Communication}. SpringerBriefs in Computer Science, 2025. \doi{10.1007/978-981-96-5132-0}.
\bibitem{liang2025} Z.~Liang, K.~Niu, J.~Xu, and P.~Zhang, ``Semantic arithmetic coding using synonymous mappings,'' \emph{Entropy}, vol.~27, no.~4, art.~429, 2025. \doi{10.3390/e27040429}.
\bibitem{xu2025} J.~Xu, T.~Du, Z.~Liang, K.~Niu, and P.~Zhang, ``Semantic Markov chain using synonymous mapping,'' \emph{Tsinghua Sci. Technol.}, vol.~32, no.~1, pp.~116--126, 2027. \doi{10.26599/tst.2025.9010064}.
\bibitem{geiger2014} B.~C.~Geiger and C.~Temmel, ``Lumpings of Markov chains, entropy rate preservation, and higher-order lumpability,'' \emph{J. Appl. Probab.}, 2014. \doi{10.1239/jap/1421763331}.
\bibitem{hinton1993} G.~E.~Hinton and D.~van~Camp, ``Keeping the neural networks simple by minimizing the description length of the weights,'' in \emph{Proc. COLT}, 1993, pp.~5--13. \doi{10.1145/168304.168306}.
\bibitem{frey1997} B.~J.~Frey, ``Efficient stochastic source coding and an application to a Bayesian network source model,'' \emph{Comput. J.}, vol.~40, no.~2--3, pp.~157--165, 1997. \doi{10.1093/comjnl/40.2_and_3.157}.
\bibitem{rissanen1979} J.~Rissanen and G.~G.~Langdon, ``Arithmetic coding,'' \emph{IBM J. Res. Dev.}, 1979. \doi{10.1147/rd.232.0149}.
\bibitem{witten1987} I.~H.~Witten, R.~M.~Neal, and J.~G.~Cleary, ``Arithmetic coding for data compression,'' \emph{Commun. ACM}, vol.~30, no.~6, pp.~520--540, 1987. \doi{10.1145/214762.214771}.
\bibitem{cover2006} T.~M.~Cover and J.~A.~Thomas, \emph{Elements of Information Theory}. Wiley, 2005. \doi{10.1002/047174882X}.
\end{thebibliography}
\end{document}
